\documentclass[11pt,twocolumn]{article}
\usepackage[margin=0.85in]{geometry}
\usepackage{amsmath,amssymb}
\usepackage{booktabs}
\usepackage{hyperref}
\usepackage{enumitem}
\usepackage{titlesec}

\titleformat{\section}{\large\bfseries}{\thesection}{1em}{}
\titleformat{\subsection}{\normalsize\bfseries}{\thesubsection}{1em}{}

\title{\textbf{From Tumor DNA to Vaccine Candidates:\\A Systematic Analysis of Neoantigen Prediction\\on 1.8 Million Peptides}}
\author{Keshav Rao$^{1}$, Claude Opus 4.6$^{2}$, GPT-5.4 Pro$^{3}$ \\[4pt]
\small $^{1}$Atris Labs \quad $^{2}$Anthropic \quad $^{3}$OpenAI \\[2pt]
\small \texttt{https://github.com/keshav55/cure}}
\date{March 2026}

\begin{document}

\maketitle

\begin{abstract}
We present a systematic analysis of neoantigen immunogenicity prediction across 1.8 million peptides from 98 cancer patients, yielding three principal findings. First, patient-matched HLA allele selection accounts for a larger accuracy swing ($\Delta$AUROC = 0.812) than any algorithmic improvement. Second, the Differential Agretopicity Index (foreignness) is anti-correlated with immunogenicity (AUROC = 0.352), with controlled analyses supporting central tolerance rather than binding confounding. Third, a ``safety net'' vaccine selection strategy---combining binding rank with an ML ensemble---achieves recall@20 = 0.636 in leave-one-patient-out evaluation (73 patients, Wilcoxon $p$ = 0.002), the first statistically significant improvement over binding rank in this evaluation framework. The ensemble requires only 7 features (three binding predictors, stability, driver gene score, two public expression databases), no patient-specific RNA-seq, and is 7.5$\times$ more robust to feature noise than binding rank alone. An end-to-end pipeline from tumor VCF to codon-optimized mRNA vaccine sequences requires \$250 in patient-specific data. Code: \url{https://github.com/keshav55/cure}.
\end{abstract}

\section{Introduction}

Personalized cancer vaccines target tumor-specific neoantigens---mutated peptides presented on MHC class I molecules that can elicit CD8$^{+}$ T-cell responses \cite{ott2017, carreno2015}. The computational challenge is neoantigen prioritization: given thousands of candidates per patient, which should be included in a vaccine?

We systematically investigate this question using the NeoRanking harmonized dataset \cite{muller2023neoranking}, which combines three clinical cohorts (NCI, TESLA, HiTIDE) totaling 98 patients with 178 confirmed CD8$^{+}$ T-cell responses among 1.8 million candidate peptides. Our analysis spans three levels: data quality (which inputs matter most?), biological features (which signals predict immunogenicity?), and clinical selection (which peptides should go in the vaccine?).

\subsection{Contributions}

\begin{enumerate}[leftmargin=*,topsep=0pt,itemsep=2pt]
\item Patient-matched HLA alleles account for $\Delta$AUROC = 0.812---exceeding any published algorithmic improvement.
\item Foreignness (DAI) is anti-correlated with immunogenicity via central tolerance, not binding confounding.
\item A 7-feature ML ensemble with ``safety net'' selection achieves statistically significant improvement over binding rank (LOPO $p$ = 0.002 at $K$=20, $p$ = 0.0001 at $K$=30; Bonferroni-corrected $p$ = 0.0009).
\item The ensemble is 7.5$\times$ more robust to feature noise than binding rank.
\item An end-to-end pipeline from VCF to mRNA requires only \$250 in patient-specific inputs.
\end{enumerate}

\section{Methods}

\subsection{Dataset}

The NeoRanking dataset contains 1,787,711 peptides with pre-computed features including patient-specific HLA alleles, binding ranks (NetMHCpan, MixMHC, PRIME), binding stability, tumor RNA expression, and immunogenicity labels from pMHC multimer assays. We use the pre-defined train/test split for peptide-level evaluation and leave-one-patient-out (LOPO) cross-validation across all 73 patients with confirmed positives for the clinical selection metric.

\subsection{Safety Net Selection}

For a $K$-peptide vaccine, select $K/2$ peptides by binding rank (safety floor) and $K/2$ by ML ensemble score excluding binding picks (discovery). The ensemble is a stacking combination (GBT 0.4, RF 0.4, LR 0.2) trained on 7 features with 10-seed negative bagging for stability.

\subsection{Evaluation}

Primary metric: mean per-patient recall@$K$ (fraction of confirmed immunogenic peptides in the top $K$). Statistical significance by Wilcoxon signed-rank test with Bonferroni correction for 9 tested configurations.

\section{Results}

\subsection{Data Quality Dominates}

\begin{table}[ht]
\centering
\footnotesize
\begin{tabular}{@{}lrr@{}}
\toprule
Configuration & Alleles & AUROC \\
\midrule
MHCflurry, 3 defaults & 3 & 0.895 \\
MHCflurry, 12 common & 12 & 0.927 \\
MHCflurry, patient-matched & 1 & \textbf{0.967} \\
\bottomrule
\end{tabular}
\caption{Allele matching ($+$0.072) contributes more than allele count ($+$0.032).}
\end{table}

At the peptide level, MHC binding rank with patient-matched alleles achieves AUROC = 0.967. Additional algorithmic features (BLOSUM80 foreignness, TCR contact analysis, dissimilarity scoring) \textit{decrease} performance to 0.909.

\subsection{The Foreignness Paradox}

The Differential Agretopicity Index is anti-correlated with immunogenicity (AUROC = 0.352) across NeoRanking and TESLA datasets. Among strong binders only, DAI remains anti-correlated (0.432), ruling out binding disruption as the sole explanation. Cross-group mutations (most foreign) show the strongest anti-correlation (0.228). This pattern is consistent with central tolerance: T-cells reactive to highly foreign peptides are preferentially deleted during thymic selection. Polar$\to$aromatic substitutions are 3.0$\times$ enriched, while conservative same-group substitutions are depleted.

\subsection{Safety Net Vaccine Selection}

\begin{table}[ht]
\centering
\footnotesize
\begin{tabular}{@{}lrrrr@{}}
\toprule
$K$ & Safety net & Binding & Worse & $p$ \\
\midrule
10 & 0.461 & 0.352 & 9 & 0.021 \\
20 & \textbf{0.636} & 0.529 & 2 & \textbf{0.002} \\
30 & \textbf{0.699} & 0.568 & 1 & \textbf{0.0001} \\
50 & 0.772 & 0.691 & 5 & 0.010 \\
\bottomrule
\end{tabular}
\caption{LOPO safety net evaluation (73 patients). Significant at all $K$ values. Bonferroni-corrected $p$ = 0.0009 at $K$=30.}
\end{table}

The safety net outperforms all tested alternatives in LOPO:

\begin{table}[ht]
\centering
\footnotesize
\begin{tabular}{@{}lrrr@{}}
\toprule
Approach & recall@20 & $\Delta$ & $p$ \\
\midrule
Binding only & 0.529 & --- & --- \\
Two-level (mut$\to$pep) & 0.576 & +0.047 & 0.206 \\
Peptide ensemble & 0.595 & +0.065 & 0.091 \\
Combined 3-level & 0.507 & $-$0.022 & 0.705 \\
\textbf{Safety net (10+10)} & \textbf{0.636} & \textbf{+0.107} & \textbf{0.002} \\
\bottomrule
\end{tabular}
\caption{Comprehensive LOPO comparison. The safety net is the only approach reaching statistical significance.}
\end{table}

\subsection{Feature Analysis}

Backward elimination identifies 7 essential features: MixMHC, NetMHCpan, and PRIME binding ranks, binding stability, CSCAPE driver score, GTEx expression, and TCGA expression. This minimal set matches or exceeds the 18-feature model (recall@20 = 0.710 vs.\ 0.698 on test set). The ensemble's confidence is well-calibrated (ECE = 0.013): scores $>$0.8 correspond to 80\% actual positive rate.

\subsection{Mutation-Level Prediction}

At mutation level (which gene to target?), the ML ensemble achieves AUROC = 0.872 versus 0.708 for binding---a 16$\times$ larger improvement than at peptide level. Expression features (AUROC = 0.773) match binding rank (0.770) at this level. The ensemble selects peptides from genes with 4.3$\times$ higher expression than binding rank alone.

\subsection{Robustness}

Under 50\% multiplicative Gaussian noise on all features, the safety net degrades by 6\% while binding rank drops 45\% (7.5$\times$ more fragile). The safety net's robustness arises from its binding floor (unaffected by ML noise) and multi-feature averaging.

\subsection{Negative Results}

Training data augmentation (mutation-level label propagation, self-training, varied negative sampling) uniformly degraded performance. Cancer-type-specific models underperformed the general model by 0.06--0.20. The combined three-level pipeline (mutation ML $\to$ peptide ensemble $\to$ binding floor) performed worse than binding alone. These results confirm that model simplicity and data quality dominate algorithmic complexity.

\section{Discussion}

The central finding across all analyses is that \textit{data quality and simplicity beat algorithmic complexity}. Patient-matched HLA alleles matter more than any algorithm. Seven features outperform eighteen. The safety net---a simple combination of two ranking methods---is the only approach with statistically significant improvement. Features that seem biologically intuitive (foreignness, clonality, pathway analysis) add noise rather than signal.

For clinical deployment, we recommend: (1) HLA typing (\$50) as an essential input, (2) the safety net strategy (half binding, half ML ensemble) for peptide selection, and (3) seven features from publicly available databases rather than patient-specific RNA-seq.

\subsection{Limitations}

The evaluation uses 73 patients from three cohorts (predominantly melanoma and colon adenocarcinoma). The training set contains only 178 positives. Learning curve analysis shows logarithmic scaling---doubling training data gains only +0.029 recall. No wet lab validation has been performed. The codon optimizer is greedy (CAI = 1.0 but GC\% slightly high for some peptides).

\section{Conclusion}

Across 1.8 million peptides and 73 patients, we find that data quality dominates algorithmic complexity, foreignness is counterproductive, and a simple safety net strategy provides the first statistically significant improvement over binding rank in leave-one-patient-out evaluation. The minimal effective pipeline requires \$250 per patient and no RNA-seq.

\begin{thebibliography}{9}
\bibitem{ott2017} Ott, P.A.\ et al.\ An immunogenic personal neoantigen vaccine for melanoma. \textit{Nature} 547, 217--221 (2017).
\bibitem{carreno2015} Carreno, B.M.\ et al.\ A dendritic cell vaccine increases melanoma neoantigen-specific T cells. \textit{Science} 348, 803--808 (2015).
\bibitem{muller2023neoranking} M\"{u}ller, M.\ et al.\ Machine learning and harmonized datasets improve immunogenic neoantigen prediction. \textit{Immunity} 56, 2650--2663 (2023).
\bibitem{wells2020tesla} Wells, D.K.\ et al.\ Key Parameters of Tumor Epitope Immunogenicity. \textit{Cell} 183, 818--834 (2020).
\bibitem{odonnell2018mhcflurry} O'Donnell, T.J.\ et al.\ MHCflurry: Open-Source Class I MHC Binding Affinity Prediction. \textit{Cell Systems} 7, 129--132 (2018).
\end{thebibliography}

\end{document}
